% ============================================
% MEM-04 (EN)
% English translation — SPEC-935-R355
% ============================================
\section*{\textbf{MEM-04} --- The Trench}
\markboth{MEM-04 --- The Trench}{}

\noindent \textit{The encounter}

\noindent \textbf{1915 --- On the road from Senador Pompeu to Fortaleza}

\vspace{0.5em}\noindent\rule{\textwidth}{0.5pt}\vspace{0.5em}

\noindent My father died in the morning. The sun rose and he was already stiff, his eyes open to a cloudless sky. My mother passed her hand over his eyes and said: \textit{\textquotedblleft{}Come, Joaquim. Your father is gone.\textquotedblright{}}

\noindent I was eight years old.

\noindent We had been walking for four days. Before, we were six. Now we were three: me, my mother, and old Anselmo, who was no relative but walked with us because alone no one survives. Old Anselmo said he had seen the lagoon dry up in 1877. Said that was a real drought. I looked at him and saw that he did not believe what he was saying. His eyes were just like my father's.

\noindent On the fifth day, we met the convoy.

\noindent They were soldiers. Thin, dirty, in patched uniforms. One carried a rifle. The other carried a whip.

\noindent \textit{\textquotedblleft{}Where are you going?\textquotedblright{}}, asked the one with the whip.

\noindent My mother answered that we were going to the coast. That she had heard there was food there.

\noindent The soldier laughed. \textit{\textquotedblleft{}Food? The coast is full of retirantes just like you. Everybody is dying on the beach.\textquotedblright{}}

\noindent My mother did not answer.

\noindent \textit{\textquotedblleft{}Come with us\textquotedblright{}\textit{, said the soldier. }\textquotedblleft{}There's a camp at Alagadiço, near Fortaleza. They give food there.\textquotedblright{}}

\noindent She hesitated. I felt her hand tighten on mine.

\noindent \textit{\textquotedblleft{}It's an order\textquotedblright{}}, said the soldier, and the whip cracked in the air. It hit no one. But the sound stayed. As if it had marked us.

\vspace{0.5em}\noindent\rule{\textwidth}{0.5pt}\vspace{0.5em}

\noindent The Government Pen was a barbed-wire fence around a clearing. No trees. No shade. People. Lots of people. More people than I had seen in my whole life.

\noindent They lay on the ground, leaning on one another. Some moved their mouths without a sound. Others did not move at all.

\noindent \textit{\textquotedblleft{}Molambudos\textquotedblright{}}, someone said.

\noindent I did not know what it meant. Later I learned. It means rags. Human rags. People who are no longer people because the drought took everything they had inside.

\noindent We stayed. My mother found a spot near the fence. Old Anselmo sat down and never got up again. On the second day, he said: \textit{\textquotedblleft{}I'm going.\textquotedblright{}} And he went.

\noindent My mother began to cough.

\vspace{0.5em}\noindent\rule{\textwidth}{0.5pt}\vspace{0.5em}

\noindent The trench lay behind the pen, about two hundred yards away. Every night the soldiers carried the bodies there. During the day you could see the vultures. They made a circle in the sky, turning, turning, as if the sky were a giant eye.

\noindent I went to the trench one night when my mother could not stop coughing.

\noindent I did not go out of curiosity.
I went because I heard something.

\noindent It was a voice.

\noindent Not a human voice --- but a voice. It was singing. Not a song with words. A melody. Low. Hoarse. A melody that rose and fell like wind in an empty place.

\noindent It sang:

\noindent \textit{\textquotedblleft{}Whoever is hungry sits. Whoever is hungry waits. Whoever is hungry eats whatever comes.\textquotedblright{}}

\noindent I stopped at the edge of the trench.

\noindent Down below, among the bodies, something moved.

\noindent I should have run. I should have gone back to the pen. I should have buried my face in my mother's skirt and forgotten that I had heard anything.

\noindent But I was a child.
And children who go hungry do not flee from food.
They investigate.

\vspace{0.5em}\noindent\rule{\textwidth}{0.5pt}\vspace{0.5em}

\noindent The trench was not deep. About ten feet. The bodies piled up like sacks of flour. Some were naked. Others had their hands tied. Their eyes were open. All of their eyes were open.

\noindent In the middle of the pile, something moved carefully.

\noindent It was not the movement of someone struggling to get out.
It was the movement of someone \textit{arranging} the bodies around herself.

\noindent An arm emerged. Not a human arm --- it was too thin, too long, with joints where joints should not be. The hand had too many fingers, and each finger had a black nail that shone like wet glass.

\noindent The hand took a loose arm from the pile and set it aside. Carefully. Like someone setting the table before a meal.

\noindent Another arm emerged. Then the head.

\noindent The head was almost human. Almost. The shape was that of a skull, but the skin still covered the bone --- thin, gray skin, stretched like wet paper over the bones of the face.

\noindent The nose was a hole.
The mouth was a line.
The eyes were two slits.

\noindent And inside the slits, a yellow glow.

\noindent She looked at me.

\noindent She did not move. She did nothing. She only looked at me, head tilted, like a dog hearing a strange sound.

\noindent I did not run.

\noindent She began to climb.

\noindent Not the way a person climbs. The way a spider does. Her arms and legs moved in coordination, finding support on the bodies, sliding among them without haste. Bones crackled under her weight. The flesh gave way.

\noindent She came out of the trench.

\noindent She stood in front of me.

\noindent She was immense --- not tall, but \textit{wide}. Her body was a cage of ribs turned outward, as if the skeleton had grown inside out. Between the ribs, something pulsed. A large, dark organ, covered with veins that glistened like worms.

\noindent Her smell reached me.

\noindent It was not the smell of death. It was the smell of \textit{food}. Roasted meat. Thick broth. The smell that came from Dona Sinhá's kitchen before the drought. The smell I had not felt in months.

\noindent My stomach ached.

\noindent \textit{\textquotedblleft{}Are you hungry?\textquotedblright{}}

\noindent Her voice was soft. Almost sweet. It did not match her body. It was the voice of a grandmother asking if her grandson had eaten.

\noindent I did not answer.

\noindent \textit{\textquotedblleft{}You are.\textquotedblright{}\textit{ She tilted her head to the other side. }\textquotedblleft{}I can feel it. Your hunger is great. As great as mine.\textquotedblright{}}

\noindent \textit{\textquotedblleft{}Ma'am...\textquotedblright{}\textit{ My voice came out weak. }\textquotedblleft{}What are you?\textquotedblright{}}

\noindent She thought. The organ between her ribs pulsed harder.

\noindent \textit{\textquotedblleft{}I am what is left\textquotedblright{}\textit{, she said. }\textquotedblleft{}I am what grows where much dies. I am what the earth cannot swallow.\textquotedblright{}}

\noindent \textit{\textquotedblleft{}Do you eat the dead?\textquotedblright{}}

\noindent \textit{\textquotedblleft{}Someone has to eat.\textquotedblright{}}

\noindent She said it with naturalness. Like someone saying \textquotedblleft{}the sun rises every day\textquotedblright{}.

\noindent \textit{\textquotedblleft{}The living don't want to. The living bury. The living cry. Then they go away. The dead stay. Someone has to make use of them.\textquotedblright{}}

\noindent She reached her hand out to me. The big hand, with long fingers, with black nails.

\noindent \textit{\textquotedblleft{}Do you want some?\textquotedblright{}}

\noindent \textit{\textquotedblleft{}Want what?\textquotedblright{}}

\noindent \textit{\textquotedblleft{}To eat.\textquotedblright{}}

\noindent I looked at her hand. I looked at the trench. I looked at the bodies.

\noindent \textit{\textquotedblleft{}No\textquotedblright{}\textit{, I said. }\textquotedblleft{}I don't want to eat them.\textquotedblright{}}

\noindent She laughed. The laugh sounded like dry leaves stirring.

\noindent \textit{\textquotedblleft{}It's not them I'm offering, boy.\textquotedblright{}}

\noindent \textit{\textquotedblleft{}Then what?\textquotedblright{}}

\noindent She crouched down. Her face came level with mine. The smell of food grew stronger. My stomach ached again.

\noindent \textit{\textquotedblleft{}I eat for you\textquotedblright{}\textit{, she said. }\textquotedblleft{}I eat every day. I eat those who die. I eat those no one wanted. And you will never feel hunger again.\textquotedblright{}}

\noindent She waited.

\noindent \textit{\textquotedblleft{}Only...\textquotedblright{}}

\noindent \textit{\textquotedblleft{}Only what?\textquotedblright{}}

\noindent She smiled. Her mouth opened more than it should. The teeth were in too many rows. Rows that could not fit in a human mouth.

\noindent \textit{\textquotedblleft{}The hunger doesn't go away. It stays. Only it's no longer yours. It's mine. And I will live inside you to feel it.\textquotedblright{}}

\noindent She touched my chest with one finger.

\noindent \textit{\textquotedblleft{}Do you accept?\textquotedblright{}}

\noindent I should have said no.

\noindent I should have spat on her hand. I should have run. I should have screamed.

\noindent But I was eight years old.
And I was so hungry.
So hungry that the no would not come out.

\noindent \textit{\textquotedblleft{}Yes\textquotedblright{}}, I said.

\noindent She opened her mouth.

\noindent It was not a kiss. It was not a bite. It was a breath. A warm breath that came out of her mouth and entered mine. It went down my throat. It filled my chest. It drained into my stomach.

\noindent She was inside me.

\noindent \textit{\textquotedblleft{}Done\textquotedblright{}\textit{, she said. }\textquotedblleft{}Now the hunger is mine.\textquotedblright{}}

\noindent She stood up. Looked at the sky. The moon was shining.

\noindent \textit{\textquotedblleft{}Go\textquotedblright{}\textit{, she said. }\textquotedblleft{}Go live. We'll meet again later.\textquotedblright{}}

\noindent She went back to the trench. She climbed down among the bodies. She lay down. And closed her eyes.

\noindent I stood at the edge, watching.

\noindent I no longer felt hunger.
But I felt a weight.
A weight at the bottom of my stomach.
A weight that breathed.\NE{The passage of the pact at the trench is the only one in the entire diary in which Joaquim writes with his left hand. The examiners who studied the manuscript noted: the handwriting changes abruptly at the line \textquotedblleft{}Yes, I said\textquotedblright{}, the pen pressure on the paper is different, and the stroke of the letters veers to the left --- as if another hand were guiding his. The editor did not intervene. The choice to preserve the original handwriting is an ethical and documentary choice.}

\vspace{0.5em}\noindent\rule{\textwidth}{0.5pt}\vspace{0.5em}

\noindent That night, at the edge of a trench of the dead, in the middle of a concentration camp, an eight-year-old boy made a deal with a monster.

\noindent The monster ate his hunger.

\noindent But it left the hole.

\noindent The hole stayed.

\noindent And everything that entered the hole afterward --- years, blood, people, memories --- never filled it.

\noindent Because the hole is not mine.

\noindent It is hers.

\noindent She has lived inside me since that night.

\noindent She eats. I feel it.

\noindent She sleeps. I wake.

\noindent She waits. I live.

\noindent She is waiting for you too.

\vspace{0.5em}
\noindent\textit{\textbf{↪ Routes:} \rota{MEM-05} \rotasep \rota{DOC-01} \rotasep \rota{LUC-01}}
